\documentclass[10pt]{article}
\usepackage[margin=20mm]{geometry}
\usepackage{amsmath,amssymb,mathtools,fontspec,unicode-math,xurl,hyperref}
\IfFontExistsTF{DejaVu Serif}{\setmainfont{DejaVu Serif}}{\setmainfont{Latin Modern Roman}}
\setmathfont{latinmodern-math.otf}
\setlength{\parindent}{0pt}\setlength{\parskip}{6pt}\sloppy
\title{Initial technical review of
Boundary Harnack principle for nonlocal operators on metric measure spaces}
\author{AI-assisted technical review}
\date{}
\begin{document}
\maketitle
Abstract: This review identifies a missing exclusion of driftless compound Poisson subordinators in the stated scope of Theorem 1.6, and a separate, locally repairable coordinate-maximization error in the appendix. An explicit exponential jump law satisfies the stated density conditions but produces an atom in every resolvent; the Green density required by the paper’s definition then does not exist. This invalidates the unrestricted application of Corollary 1.5, without disproving a boundary Harnack inequality or the later theorem with a positive Gaussian component. The appendix error has an explicit repair preserving the claimed estimate. The reviewed object is arXiv:2312.02455v1, not a verified copy of the journal version.\par
\section*{Identity, scope, and conclusion}
The assigned paper is by Zhen-Qing Chen and Jie-Ming Wang, DOI 10.1214/24-AOP1734 (https://doi.org/10.1214/24-AOP1734). The acquired source is the 49-page arXiv version dated 5 December 2023, [CW]; its title page uses the spelling non-local. All page locators below refer to this exact PDF, whose physical and printed page numbers agree. The source SHA256 is\par
> 00000TOKEN.\par
The principal dependency examined is \[\begin{gathered}
\text{(A1)--(A4) and Green comparison}\ \longrightarrow\ \text{Theorem 1.3},\\
\text{Theorem 1.3}\ \longrightarrow\ \text{Corollary 1.5}
\ \longrightarrow\ \text{Theorem 1.6}.
\end{gathered}\] The later geometric result, Theorem 1.10, additionally assumes positive subordinator drift. Its critical-cone mechanism was examined separately. The appendix estimate behind Theorems 2.8 and 2.9 was also checked.\par
Text on pages 1–20, 23–31, and 38–46 was read. Reading these pages does not certify every proof on them. The decisive checks concern the hypotheses and occupation kernel in Problem 1, the coordinate estimate in Problem 2, and the two-dimensional critical-cone calculation in Section 4. Pages 21–22, 32–37, and 47–49 and the full proofs in external references have not received complete scientific review. Images of all source pages used as decisive evidence were inspected; the separate coverage record lists them. No author code was executed.\par
There are two established problems: one Major problem in the scope and deduction of Theorem 1.6, and one Minor proof error in Lemma 4.1. Neither finding establishes that the paper’s principal positive-Gaussian boundary Harnack theorem is false. This is an initial technical report; scientific acceptance remains pending independent review.\par
\section*{A missing hypothesis in the subordinate Brownian application}
Problem 1 (Major: the stated class includes processes without the required Green density).  On page 8, the paragraph preceding (1.11) asserts transition-density regularity for any subordinate Brownian motion. Theorem 1.6 on the same page assumes a discontinuous subordinator with a Lévy density satisfying (1.12), but does not require positive drift or infinite jump activity. These conditions permit a process violating (A1), which is explicitly required in Corollary 1.5 on page 7. For that process the Green function used in (1.6) cannot satisfy Definition 1.1 on page 5.\par
\subsection*{Exact hypotheses and an admissible example}
Equation reference: eq:assumptions. The subordinator exponent in (1.11) is \[\phi(\lambda)=b\lambda+\int_0^\infty(1-e^{-\lambda s})\mu(s)\,ds,
 \qquad b\geq0,\] where \(\mu\) is nonzero and \(\int_0^\infty(1\wedge s)\mu(s)\,ds<\infty\). The two displayed assumptions (1.12) are, for a fixed finite constant \(c\), \[
 \mu(s)\leq c\mu(2s)\quad(0<s<8),\qquad
 \mu(s)\leq c\mu(s+1)\quad(s>1).\] Take \(b=0\) and \(\mu(s)=e^{-s}\), \(s>0\). Then \[\int_0^\infty\mu(s)\,ds=1,\qquad
 \int_0^\infty(1\wedge s)\mu(s)\,ds=1-e^{-1},\qquad
 \phi(\lambda)=\frac{\lambda}{1+\lambda}.\] Both inequalities in [eq:assumptions] (\#eq:assumptions) hold with \(c=e^8\), since their ratios are respectively \(e^s\) and \(e\). This is a discontinuous compound Poisson subordinator with rate one and exponential jumps. Its density is completely monotone, so restricting the example to complete Bernstein exponents does not remove it.\par
Let \(W\) be the independent Brownian motion used in the paper and set \(X_t=W_{S_t}\) in \(\mathbb R^d\), for any \(d\geq1\). The process is a symmetric compound Poisson process; its jump distribution is a mixture of nondegenerate Gaussian distributions and is absolutely continuous. The holding component is the obstruction.\par
Equation reference: eq:atom. Claim 1.  For every \(x\in\mathbb R^d\), \(t>0\), and \(\alpha>0\), \[
 \mathbb P_x(X_t=x)=e^{-t},
 \qquad G^\alpha(x,\{x\})=\frac{1}{\alpha+1}.\] For every open set \(U\) containing \(x\), the killed occupation measure has mass one at \(\{x\}\).\par
Proof. There is no subordinator jump before \(t\) with probability \(e^{-t}\). Conditional on \(S_t>0\), \(W_{S_t}\) has a Gaussian-mixture density and therefore no atom at \(x\). Integrating the first equality in [eq:atom] (\#eq:atom) against \(e^{-\alpha t}\,dt\) gives the second.\par
Let \(T_1\) be the first jump time. It has mean one, and the process stays at \(x\) up to \(T_1\), before exiting any open \(U\) containing \(x\). The locations after each positive number of jumps have continuous distributions; with probability one none of these countably many locations equals \(x\). Consequently \[\mathbb E_x\int_0^{\tau_U}\mathbf1_{\{x\}}(X_s)\,ds
   =\mathbb ET_1=1.\] An occupation density with respect to Lebesgue measure would give zero for this bounded Borel test function. This contradicts the defining identity in Definition 1.1. ◻\par
\subsection*{Dependency and impact}
Assumption (A1), page 4, requires absolute continuity of every positive resolvent for both dual processes. Equation [eq:atom] (\#eq:atom) violates it. Theorem 1.3 explicitly assumes (A1), and Corollary 1.5 carries it into the Euclidean Lévy-process application. Page 8 supplies the density assertion to justify the next application. That assertion is false, and the conditions printed in Theorem 1.6 do not supply its missing hypothesis.\par
This is more than the failure of an optional regularity remark: Theorem 1.6 states an equivalence with a comparison of Green functions defined by an occupation-density identity. Choose, for example, a bounded open ball \(D\); its complement has nonempty interior, so no boundary-domain restriction removes the example. For the compound Poisson process, the corresponding killed Green density does not exist in the paper’s stated sense. Thus the advertised equivalence is not well defined over its entire stated class, and its deduction from Corollary 1.5 is invalid for that class. This is the reason for the Major classification.\par
\subsection*{Alternative readings, counterevidence, and repair limits}
The word discontinuous permits compound Poisson subordinators; it does not impose infinite activity. Formula (1.11) expressly permits \(b=0\). The positive-drift convention introduced at the start of Section 3, page 23, belongs to that later section. Moreover, page 29, after (3.18), expressly observes that the jump-kernel comparison does not require positive drift. That observation establishes (A3)’ for a wider class; it does not establish (A1).\par
One could study an off-diagonal density after removing the holding atom. It would not satisfy Definition 1.1 for all bounded nonnegative Borel functions, as the singleton test above shows. Such a formulation would need a changed definition and a new check of the theorem’s proof.\par
Adding \(b>0\), or excluding finite-activity driftless subordinators by requiring \(\mu(0,\infty)=\infty\), removes this particular resolvent obstruction: \(S_t>0\) almost surely for \(t>0\), and Gaussian mixing then gives absolute continuity. This observation alone does not verify all remaining assumptions of a revised theorem, and does not assert continuity of an arbitrary resulting density. An explicit (A1) hypothesis is another direct way to state the required framework.\par
There is no counterexample here to the boundary Harnack inequality itself. The abstract metric-space theorem already assumes (A1), and Theorem 1.10 on page 10 expressly requires \(b>0\); neither is refuted by this example. The symbolic integration and atom calculation were independently reproduced in the saved reviewer script `verification/0512/evidence/independent\_checks.py`; the analytic argument above is the decisive evidence.\par
\section*{An invalid coordinate aggregation in the appendix}
Problem 2 (Minor: two coordinate maxima cannot be separated).  In Lemma 4.1, Case (i), page 45, the displayed chain immediately after (4.5) replaces the maximum of a sum of two coordinate-dependent terms by the sum of their separate maxima. This equality is false for admissible Lévy triples. Keeping the joint maximum repairs the estimate with the same stated order.\par
\subsection*{Hypotheses, exact step, and counterexample}
Equation reference: eq:falsemax. The appendix treats any Lévy process in \(\mathbb R^d\) with covariance matrix \(A\), drift vector, and Lévy measure \(\nu\). For fixed \(r>0\), write \[B_i=\int_{|z|\leq2r}z_i^2\,\nu(dz),\qquad
 M=\max_{1\leq i\leq d}\left(\tfrac12a_{ii}+c_2 B_i\right),\] where \(c_2>0\) is the constant in (4.5). The infimum of the exponential integrands in that display is \(\exp(-t\omega^2 M)\). The next equality instead uses \[
 \tfrac12\max_i a_{ii}+c_2\max_i B_i .\] The maximizing indices need not coincide.\par
For a numerical instance satisfying all hypotheses of Case (i), take \(d=2\), \(r=1\), zero drift, \[A=\operatorname{diag}(1/2,0),\qquad
 \nu=\tfrac12\delta_{(0,1)}+\tfrac12\delta_{(0,-1)},\qquad c_2=1/4.\] The constant \(1/4\) is valid in the cosine estimate for \(|\omega z_i|\leq1\), since \(1-\cos u\geq u^2/4\) there. The measure is a finite symmetric Lévy measure, \(B_1=0\), \(B_2=1\), and \[K(2r)=3/8,\qquad G(2r)=L(2r)=0.\] Thus \(K(2r)\geq\max\{G(2r),L(2r)\}\), exactly the case under discussion. But \(M=1/4\), whereas [eq:falsemax] (\#eq:falsemax) equals \(1/2\). For every \(t\omega^2>0\) the printed equality therefore changes the exponential to a strictly smaller value without justification. More generally, for any chosen \(c_2>0\), taking \(A=\operatorname{diag}(2c_2,0)\) with the same measure gives the same factor-two discrepancy.\par
\subsection*{Repair and dependency}
The required bound follows without separating the maxima: \[\begin{aligned}
 M&\geq\frac1d\sum_{i=1}^d
       \left(\tfrac12 a_{ii}+c_2 B_i\right) \notag\\
  &\geq \frac{\min\{1/2,c_2\}}{d}
     \left(\operatorname{tr}A+\sum_i B_i\right)
   =\frac{\min\{1/2,c_2\}}d\,4r^2K(2r).\end{aligned}\] Inserting this in the preceding Gaussian integral yields \(C/\sqrt{tK(2r)}\), precisely the estimate needed in Case (i). The subsequent use of independent increments then gives its desired \(C/[t\Phi(2r)]\) bound. Thus this is a demonstrated invalid algebraic step, with a complete local repair. It does not refute Lemma 4.1, Pruitt’s exit-time estimate in Theorem 2.8, or the application to (A4) in Theorem 2.9. Its limited impact supports Minor severity. The exact rational example and the case hypothesis were independently checked in `verification/0512/evidence/coordinate\_max\_check.py`.\par
\section*{A positive boundary-case check}
Equation reference: eq:sector. The critical angle in dimension two is \(\pi/4\), corresponding to a sector of opening \(\pi/2\). Lemma 3.14, pages 39–40, predicts the additional logarithm in the Brownian mean exit time at this angle. There is a direct independent calculation for Brownian generator \(\Delta\) on the unit quarter-disk. Write polar coordinates \(0<\theta<\pi/2\). The solution of \(\Delta u=-1\), with zero boundary values, is \[
 u(r,\theta)=\frac{r^2\log(1/r)}{\pi}\sin(2\theta)
 +\sum_{\substack{k\geq3\\k\ {\rm odd}}}
 \frac{r^2-r^{2k}}{\pi k(k^2-1)}\sin(2k\theta).\] Indeed, the sine expansion of the constant function on this interval has coefficients \(4/(\pi k)\) for odd \(k\). The first radial operator sends \(r^2\log(1/r)\) to \(-4\); the higher modes in [eq:sector] (\#eq:sector) give the remaining coefficients of \(-1\). The boundary values vanish, and the remainder on the central ray is \(O(r^2)\) by absolute summability. Consequently \[\frac{u(r,\pi/4)}{r^2\log(1/r)}\longrightarrow\frac1\pi
 \quad(r\downarrow0).\] Independent 80-digit partial-sum calculations for \(r=10^{-1},10^{-2},
10^{-3},10^{-4}\) gave approximately \(0.3134170,0.3158632,0.3166787,
0.3170865\), consistent with \(1/\pi\). The proof is the mode calculation, not a numerical extrapolation. This supports the critical logarithmic mechanism used in the later contradiction argument, but does not certify the full cone proof or its external inputs.\par
\section*{Version checks and remaining limits}
The arXiv record inspected during this review listed v1 as the only version. The author’s publication page linked that preprint. Bounded title/DOI searches did not identify a target-specific author correction; this is not a claim that no correction exists. Access to the publisher page reached an additional security check, and no publisher full-text payload was acquired or compared. The two findings therefore apply only to the exact acquired v1. They must not be attributed automatically to the 2025 journal text.\par
The report separates a false density assertion and invalid theorem application from a false boundary Harnack conclusion, and separates the appendix’s false intermediate equality from its repairable final estimate. Other proofs, versions, and external results outside the declared coverage remain unverified. The supporting records retain source and manifest bindings, reviewer calculations, actual image inspection facts, and the exact build dependencies for this report.\par
9 Z.-Q. Chen and J.-M. Wang, Boundary Harnack principle for non-local operators on metric measure spaces, arXiv:2312.02455v1, 5 December 2023, 49 pp. https://arxiv.org/abs/2312.02455. Assigned journal DOI: https://doi.org/10.1214/24-AOP1734. Z.-Q. Chen, publication page, accessed during this review on 29 October 2026. https://sites.math.washington.edu/\textasciitilde{}zchen/paper.html. Used only for version/correction discovery, not as proof evidence.\par
\end{document}
